\documentclass[10pt,a4paper]{amsart}
\usepackage[margin=19mm]{geometry}
\usepackage{amsmath,amssymb,mathtools}
\usepackage[scr=rsfs,cal=euler]{mathalfa}
\usepackage{xfrac}
\usepackage[unicode,hidelinks,bookmarks=false]{hyperref}
\newcommand{\ie}{i.e.\ }
\newcommand{\cf}{cf.\ }
\newcommand{\resp}{respectively\ }
\newcommand{\Subsection}{\subsection}
\providecommand{\nobreakdash}{\nobreak\hbox{-}}
\setlength{\emergencystretch}{2em}
\multlinegap0pt
\usepackage{bm}
\usepackage{bbm}
\usepackage{dsfont}
\usepackage{xspace}
\usepackage{cancel}
\usepackage{mathtools}
\usepackage{amsthm}
\makeatletter
\@ifundefined{theorem}{\newtheorem{theorem}{Theorem}}{}
\@ifundefined{claim}{\newtheorem{claim}[theorem]{Claim}}{}
\@ifundefined{proposition}{\newtheorem{proposition}[theorem]{Proposition}}{}
\makeatother
\newcommand{\p}{\partial}
\newcommand{\qq}{{\sf q}}   
\title[Sublinear Morse Geodesics and First Passage Percolation]{Sublinear Morse Geodesics and First Passage Percolation}
\author{Sagnik Jana, Yulan Qing}
\date{Source: arXiv:2507.07859v1 (CC BY 4.0)}
\begin{document}
\maketitle

\noindent\emph{Source and licence.} Sublinear Morse Geodesics and First Passage Percolation. Version arXiv:2507.07859v1 (CC BY 4.0), distributed under the Creative Commons Attribution 4.0 International licence, as recorded in the arXiv licence metadata for this exact version and shown on the abstract page https://arxiv.org/abs/2507.07859v1. Full rights notice: licenses/arxiv-2507.07859-CC-BY-4.0.txt.\par\smallskip

\noindent\textit{Editorial scope.} Counted unit: the Proposition whose source label is \texttt{None} in the article (source0.tex lines 1181--1189, source label \texttt{None}), delivered together with the article's own complete original proof of it (source0.tex lines 1191--1236, one \texttt{proof} environment of 46 lines). No mathematics is written, repaired or re-derived: statement and proof are the article's own text line for line. Required context retained uncounted: none. Every definition, display and label of those lines is retained unchanged. Omitted from this package and not redistributed: 1--1180, 1190--1190, 1237--2055. Nothing inside a retained excerpt is omitted or altered: every retained body is delivered exactly as the article's lines are written, so no omission receipt of any kind accompanies this package. Citations: the retained excerpts carry no citation command at all, so no bibliography entry of the article is transcribed and no reference is redistributed. Reference adaptation: none was needed and none was made; every reference inside the delivered unit resolves inside it, so no unresolved reference or citation is produced. Printed slips kept literal: no printed word, symbol, hypothesis or number of the counted statement or of its proof was altered. Layout and class: the article's own class is \texttt{amsart} with packages including inputenc, geometry, stmaryrd, amsmath, amssymb, amsthm, mathrsfs, mathtools, bbold, mathbbol, amscd, comment, graphicx, epstopdf; the wrapper is the lane's pinned \texttt{amsart} editorial wrapper at 10pt on A4, which restates the theorem-like environments the retained text needs and adds the article's own macro definitions that the retained text needs (\texttt{\textbackslash p}, \texttt{\textbackslash qq}), with extra wrapper packages (bm, bbm, dsfont, xspace, cancel, mathtools). The delivered numbering is the unit's own. Assets: no retained span contains a figure environment, a graphics-inclusion command, a PDF-page inclusion, a table or a TikZ picture; no image file is copied into this package, the package declares no asset dependency and no image file is redistributed with it. Licence: version arXiv:2507.07859v1 of this article is distributed under the Creative Commons Attribution 4.0 International licence, as recorded in the arXiv licence metadata for this exact version and shown on the abstract page https://arxiv.org/abs/2507.07859v1. A full rights notice is delivered at licenses/arxiv-2507.07859-CC-BY-4.0.txt. Characters: every retained excerpt is pure ASCII, so the delivered engine, XeLaTeX with its default Latin Modern text font, reports no missing character.
\begin{proposition}
Let $\beta$ be a bi-infinite geodesic line. Assume that  $\beta$ is $\kappa$-Morse with Morse gauge $M$. Let $p$ be any $\qq$-quasi-geodesic segments with following properties:
\begin{itemize}
\item  the length of $L(\p)\leq Cd(a,b)$
\item ${\kappa(|b|)< \frac{e|b|}{6M(C,0)[2+M(C,0)A]}}$.
\item $d(a,b)\geq e|b|$ for some constant $A$.
\end{itemize}
Then for every $C\geq 1$ and $e > 0$ there exists $D=D(C)$ such all such segments $p$ intersect the $D\cdot\kappa$-neighborhood of the middle third of the segment $[a, b]_\beta$
\end{proposition}

\begin{proof}
    Suppose $\beta$ is a $\kappa$-Morse geodesic ray starting at $\mathfrak{o}$ with Morse gauge function $M:\mathbb{R}^2 \rightarrow \mathbb{R}$. For $\kappa$ being sublinear there exist a $t$ such that 
    \[M(C,0)\kappa(t)< e\cdot t.\] Now, by way of contradiction, for some fixed $C\geq 1$ and $D(C)>0$ there exist a $C$- biLipschitz path $\p:[0,L]\rightarrow X$ joining $a,b$ with $L \leq Cd(a,b)$, for which \[\p \cap N_\kappa(\frac{1}{3}\beta[r,s], D(C))= \phi.\]  Moreover, $\beta$ is $\kappa$-Morse hence, $\p \subset N_\kappa(\beta, M(C,0))$.  
    
    
    There are two cases to consider:\newline
    \textit{Case 1.} Let $\p \cap N_\kappa(\frac{1}{3}\beta[r,s], M(C,0))= \phi $, then for every $c \in \p$ there exists a point $\eta_c\in \beta $ such that $d(c,\eta_c)\leq M(C,0)\kappa(|c|)$, where $\eta_c$ is either in $\beta_l$ and $\beta_r$, where \[\beta_l=\{x \in \beta : d(x,a)\leq \frac{1}{3}d(a,b) \text{ or } x\in [\mathfrak{o},a]_{\beta}\}\]
    and \[\beta_r= \{x \in \beta : d(y,b)\leq \frac{1}{3}d(a,b)\]
     or \[ y\in [b,\infty]_{\beta} \}.\] If the whole path is in $\kappa$- neighborhood of either $\beta_l$ or $\beta_r$ then \[d(a,b) \leq M(C,0)\kappa(|a|)+M(C,0)\kappa(|b|)< e|b|\]
     which is a contradiction. Every point on $\p$ is $M(C,0)\cdot \kappa$-close to either $\beta_l$ or $\beta_r$.
\begin{claim}There exists a $\eta$ such that \[d(\eta,\beta_l)\leq M(C,0)\kappa(|\eta|) \text{ and } d(\eta,\beta_r)\leq M(C,0)\kappa(|\eta|).\] 
\end{claim}
\begin{proof} Suppose every point is either close to only $\beta_l$ or $\beta_r$. 
$H$ be the set of all closed sub arc of $\p$ which contains in $M\cdot\kappa$ neighborhood of $\beta_l$, and similarly $K$ be the set of all closed sub arc of $\p$ which is $M\cdot\kappa$ close to $\beta_r$.  So, $\p=H \cup K$ with$H \cap K = \phi$. Which is a contradiction to the fact that $\p$ is connected. Therefore, there is a point $\eta$ such that it is $M(C,0).\kappa$-close to both $\beta_l,\beta_r.$ \newline
If $x\in \beta_l$ and $y\in \beta_r$ then $d(x,y)\geq \frac{1}{3} d(a,b)\geq \frac{e|b|}{3}$. 
 From \textit{Claim 1} we conclude that, \[ \frac{e|b|}{3} \leq d(x,y) \leq d(\eta,x)+d(\eta,y)\leq 2 M(C,0) \kappa(|\eta|)\]
From the lemma 3.1 and the fact that $d(\eta,y)\leq M(C,0)\kappa(|\eta|)$, we can say there exists $D_1= 2+M(C,0)A$, where $A$ is a constant (depends on $\kappa$) such that $\kappa (|\eta|)\leq D_1 \kappa (|y|)$. 
\newline 
Since, $\beta$ is $\kappa$- Morse geodesic then we have,
\[\frac{e|b|}{3} \leq d(x,y) \leq 2M(C,0)(2+2M(C,0)A)\kappa(|b|)\] 
therefore,
\begin{equation} 
\begin{split}
 \frac{e|b|}{3}&\leq 2M(C,0)(2+2M(C,0)A)\kappa(|b|) \\\\
   {e|b|} &\leq 6M(C,0)(1+M(C,0)A)\kappa(|b|)\\\\
    \kappa(|b|)&\geq \frac{e|b|}{6M(C,0)[2+M(C,0)A]} 
\end{split}
\end{equation}
Which is a contradiction! Hence our assumption $\beta$ being $\kappa$-Morse is wrong.\end{proof}

\textit{Case 2.}  Let $\p \cap N_\kappa(\frac{1}{3}\beta[r,s], M(C,0)) \neq \phi $. By the way of contradiction we also have $\p \cap N_\kappa(\frac{1}{3}\beta[r,s], D(C))= \phi $. Therefore, $M(C,0) > D(C)$.
\newline Now consider the $\kappa$-neighborhood $N_\kappa(\beta, D(C))$. Suppose $\mathcal{O}$ be the set of all open sub arcs which are outside of $N_\kappa(\beta, D(C))$. 
 The last time, when it enters the neighborhood, if the terminal points of the arc will be close to $\beta_l$ then for some $x \in \beta_l$, 
 \begin{align*}
d(b,x)&\leq D(C)\kappa(|b|)\\ 
   \frac{e|b|}{3} & \leq D(C)\kappa(|b|)\\
   \frac{e|b|}{3}&< \frac{e|b|\cdot D(C)}{6M(C,0)[2+2M(C,0)A]}
 \end{align*}
 After simplifying above we will get $M(C,0)\cdot A < 0$ which can not be possible. \newline
 
 Therefore, there is a close sub arc (say, $\mathfrak{h}$) inside  $N_\kappa(\beta, D(C))$ whose terminal points are close to $\beta_l$ and $\beta_r$ respectively.  The goal is to use the \textit{case 1} for this close sub arc $\mathfrak{h}$ in $\mathcal{O}^c$. Similarly like case 1, $H$ be the set of all closed sub arc of $\mathfrak{h}$ which contains in $D(C)\cdot\kappa$ neighborhood of $\beta_l$, and $K$ be the set of all closed sub arc of $\mathfrak{h}$ which is $D(C)\cdot\kappa$ close to $\beta_r$. $H\cup K= \mathfrak{h}$. Hence, there is a point $\eta$ such that it is $D(C).\kappa$-close to both $\beta_l,\beta_r$ such that
 \[ d(\eta,x)\leq D(C)\kappa(|\eta|)\]
 and \[d(\eta,y)\leq D(C)\kappa(|\eta|)\]
 for some $x \in \beta_l$ and $y\in \beta_r$. Then by applying lemma 3.1 we have, \[\kappa(|\eta|)\leq (2+2D(C)A)\kappa(|x|)\leq (2+2D(C)A)\kappa(|b|).\]
Now by applying the case 1, we can conclude that\[\kappa(|b|)\geq \frac{e|b|}{6D(C)[2+2D(C)A]} \geq \frac{e|b|}{6M(0C,0)[2+M(C,0)A]}\] 
Which is a contradiction to our assumption! This complete the proof.\end{proof}

\end{document}
